\section{问题四的模型建立与求解}

\subsection{任务分区与独立资源核算模型}

问题四以问题三的联合调度方案（承接 q3\_2 均衡预设：合并路线 $16$、架次 $29$，运输 $16$＋中继 $13$，总能耗 $67.387\,$kWh）为不变基础，将 $15$ 个服务区划分为 $K=2$ 与 $K=3$ 个任务组。约束是：组批、访问顺序、运输—中继任务与通信关系保持不变，凡同一多站架次涉及的服务区须整组同划（must-link），各组独立执行、四类资源（运输无人机、共享电池、中继无人机、中继能源组件）不得跨组调配。

为保证算法比较公平，本文只在同一条 $q3\_2\to q4\_1/q4\_2$ 数据链上比较完全枚举与局部搜索。代码包中的 $p4\_\ast$ 增强方案承接的是 $p3\_\ast$ 路线链，原子块数量和通信风险定义均不同，因此不把两套结果直接放在同一张排名表中。

\textbf{must-link 收缩.} 将“必须同组”的服务区合并为原子任务块。由问题三多站架次得到两条 must-link：$\{S_{007},S_{015}\}$ 与 $\{S_{008},S_{011}\}$，其余服务区各自成块，$15$ 个服务区收缩为 $13$ 个原子块。

\textbf{独立资源核算.} 对任一分区 $\{G_1,\dots,G_K\}$，每组 $G_j$ 承担本组服务区对应的运输与中继任务，按其架次时序独立仿真，核算并发峰值所需的运输无人机数（分机型）、共享电池数、中继无人机数与中继能源组件数。定义组间工作量\textbf{均衡系数}为各组资源规模的变异系数 $\mathrm{CV}=\sigma/\mu$，资源缺口为各组需求之和超出现有库存的部分。现有库存合计 $30$ 单位：运输无人机 A/B/C $=4/2/2$，共享电池 A/B/C $=6/4/4$，中继无人机 $=2$，中继能源组件 $=6$。

\subsection{求解方法：高复杂度主解与基础解}

\textbf{高复杂度主解（完全枚举精确分区）.} 由于原子块仅 $13$ 个，可完全枚举全部划分。$K=2$ 枚举 $4095$ 个分区，$K=3$ 枚举 $261625$ 个分区，按“缺口最小 $\to$ 规模最小 $\to$ 均衡最优”的词典序打分选优，保证在当前任务块空间内的全局最优。

$K=2$ 最优分区：
\begin{itemize}[leftmargin=2em]
  \item 组 1：$\{S_{001},S_{002},S_{003},S_{009},S_{012},S_{013}\}$；
  \item 组 2：$\{S_{004},S_{005},S_{006},S_{007},S_{008},S_{010},S_{011},S_{014},S_{015}\}$；
\end{itemize}
资源总规模 $18$，均衡系数 $0.0337$，组内冗余 $13$，资源缺口 $1$（中继无人机：两组需求合计 $3>$ 库存 $2$）。

$K=3$ 最优分区：组 1 $\{S_{001},S_{002},S_{003},S_{014}\}$、组 2 $\{S_{004},S_{005},S_{006},S_{007},S_{008},S_{011},S_{015}\}$、组 3 $\{S_{009},S_{010},S_{012},S_{013}\}$；资源总规模 $21$，均衡系数 $0.3350$，组内冗余 $10$，资源缺口 $1$（中继无人机合计 $3>2$）。

\textbf{基础解（负载均衡初始化＋模拟退火局部搜索）.} 以负载均衡贪心初始化，用 $300$ 轮模拟退火（$72\%$ 单块移动、$28\%$ 两块交换）搜索同一分区空间。$K=2$ 得规模 $23$／均衡 $0.1480$／缺口 $3$；$K=3$ 得规模 $23$／均衡 $0.0299$／缺口 $4$。它不保证全局最优，但用于说明在规模扩大、完全枚举不可承受时的近似能力。

\subsection{两种分区方式比较与缺口分析}

两方法结果汇总见表~\ref{tab:q4}，分区结构与资源分布见图~\ref{fig:q4}。以全局最优（完全枚举）为准：$K=2$ 相比 $K=3$ 资源总规模更小（$18$ vs $21$）、均衡系数更优（$0.0337$ vs $0.3350$）、冗余更高（$13$ vs $10$），二者资源缺口相同（均为 $1$）。这说明在本场景下，\textbf{分成 2 组在资源规模与组间均衡上均优于 3 组}：组数增多迫使每组独立保留基础资源，抬高总规模并加剧不均衡。

\begin{table}[htbp]
\centering
\caption{两种分区方式资源配置对比}
\label{tab:q4}
\renewcommand{\arraystretch}{1.2}
\begin{tabular}{c|c|c|c|c|c}
\toprule
方法 & 组数 & 资源总规模 & 均衡系数 & 组内冗余 & 资源缺口 \\
\midrule
完全枚举 & $K{=}2$ & 18 & 0.0337 & 13 & 1 \\
完全枚举 & $K{=}3$ & 21 & 0.3350 & 10 & 1 \\
局部搜索 & $K{=}2$ & 23 & 0.1480 & 10 & 3 \\
局部搜索 & $K{=}3$ & 23 & 0.0299 & 11 & 4 \\
\bottomrule
\end{tabular}
\end{table}

\textbf{缺口及其原因.} 全局最优分区下唯一不可避免的缺口是\textbf{中继无人机}：现有库存仅 $2$ 架，而无论分 $2$ 组还是 $3$ 组，各组为保障本组通信均至少需 $1$ 架中继，需求合计达 $3$，缺口为 $1$。该缺口是“资源不得跨组调配”这一独立执行约束的\textbf{结构性代价}——若允许中继无人机跨组共享，则 $2$ 架即可满足；一旦强制分组独立，各组必须各自留冗余，缺口随组数上升而不减。据此建议：应急指挥在采用任务分区独立执行方案时，应优先增配至少 $1$ 架中继无人机，或对中继资源保留有限的跨组机动权限。

\begin{figure}[htbp]
  \centering
  \subfloat[$K{=}2$ 任务分区结构]{\includegraphics[width=0.48\textwidth]{q4_1_01.png}}
  \hfill
  \subfloat[$K{=}3$ 任务分区结构]{\includegraphics[width=0.48\textwidth]{q4_1_02.png}}
  \\
  \subfloat[各组资源规模与均衡]{\includegraphics[width=0.48\textwidth]{q4_1_03.png}}
  \hfill
  \subfloat[资源缺口分解]{\includegraphics[width=0.48\textwidth]{q4_2_01.png}}
  \caption{问题四任务分区与资源配置结果}
  \label{fig:q4}
\end{figure}

\subsection{分区权衡空间与搜索过程分析}

为揭示分区决策的全局权衡结构，图~\ref{fig:adv_q4} 从枚举解集与搜索轨迹两方面展开。图~\ref{fig:adv_q4}(a) 将全部枚举分区映射到“资源总规模—工作量均衡系数”平面并以颜色编码资源缺口：可见总规模越小的方案（左侧）资源缺口越大（偏红），二者存在明显冲突，$K{=}2$ 与 $K{=}3$ 最优解位于缺口最小（缺口 $=1$）的低规模边界。图~\ref{fig:adv_q4}(b) 对比两种分区的组间作业时长，$K{=}2$ 两组时长接近（均衡系数 $0.034$），而 $K{=}3$ 因第三组任务偏轻而拉大组间落差（$0.335$），定量印证 $K{=}2$ 在负载均衡上更优。图~\ref{fig:adv_q4}(c) 的堆叠柱显示各组能耗以运输为主、中继为辅，第三组（$K{=}3$-G3）规模最小。图~\ref{fig:adv_q4}(d) 给出局部搜索（模拟退火）的收敛曲线，两种 $K$ 均在约 $50$ 步内快速下降并稳定，验证了求解的收敛性；其收敛值高于完全枚举最优，说明本规模下完全枚举给出的精确解确属全局最优。

\begin{figure}[htbp]
  \centering
  \includegraphics[width=0.96\textwidth]{adv_q4.png}
  \caption{分区权衡空间与搜索过程}
  \label{fig:adv_q4}
\end{figure}
